\documentclass[11pt]{article}
\usepackage[a4paper,margin=1in]{geometry}
\usepackage{fontspec}
\usepackage[boldfont=false]{xeCJK}
\setCJKmainfont{FandolSong-Regular.otf}
\usepackage{amsmath}
\usepackage{amssymb}
\usepackage{array}
\usepackage{booktabs}
\usepackage{graphicx}
\usepackage{adjustbox}
\usepackage{enumitem}
\setlist[itemize]{topsep=2pt,partopsep=0pt,parsep=0pt,itemsep=2pt,leftmargin=1.4em}
\setlist[enumerate]{topsep=2pt,partopsep=0pt,parsep=0pt,itemsep=2pt,leftmargin=1.6em}
\usepackage{parskip}
\usepackage{fvextra}
\fvset{breaklines=true,breakanywhere=true,fontsize=\small}
\setlength{\parindent}{0pt}

\begin{document}

\begin{center}
  {\LARGE\bfseries Group 17}\\[0.35em]
  {\large A Level Chemistry}
\end{center}
\vspace{0.6em}

\section*{Physical properties of the halogens}

The \textbf{halogens} 卤素 are the \textbf{Group} 族 17 elements. They exist as diatomic molecules ($\text{Cl}_2$, $\text{Br}_2$, $\text{I}_2$). Going down the group:

\begin{center}
\adjustbox{max width=\linewidth}{%
\begin{tabular}{ll}
\toprule
\textbf{Element} & \textbf{Colour and state at room temperature} \\
\midrule
chlorine & pale green gas \\
bromine & red-brown liquid \\
iodine & grey-black solid (purple vapour) \\
\bottomrule
\end{tabular}}
\end{center}

The \textbf{volatility} 挥发性 (how easily a substance turns to vapour) \textbf{decreases} down the group: chlorine is a gas, but iodine is a solid. This is because the molecules get larger and have more electrons, so the \textbf{instantaneous dipole} 瞬时偶极 and \textbf{induced dipole} 诱导偶极 forces between them get stronger. Stronger forces are harder to break, so the boiling point rises and volatility falls.

The \textbf{bond energy} 键能 (bond strength) of the $\text{X}\text{–}\text{X}$ molecules generally falls from $\text{Cl}_2$ to $\text{I}_2$, because the shared electrons are further from the nuclei in the larger atoms.

\section*{Chemical properties of the halogens and hydrogen halides}

\subsection*{Halogens as oxidising agents}

Each halogen reacts by gaining one electron to form a $1-$ ion, so it acts as an \textbf{oxidising agent} 氧化剂. This power \textbf{decreases} down the group, because the larger atoms attract an extra electron less strongly. A more reactive halogen can push out a less reactive one from its salt:

\[
\text{Cl}_2 + 2\text{KBr} \rightarrow 2\text{KCl} + \text{Br}_2
\]

\subsection*{Reactions with hydrogen}

Each halogen reacts with hydrogen to form a \textbf{hydrogen halide} 卤化氢:

\[
\text{H}_2 + \text{Cl}_2 \rightarrow 2\text{HCl}
\]

The reaction gets less vigorous down the group: chlorine reacts explosively in light, bromine needs heat, and iodine reacts slowly and only partly.

\subsection*{Thermal stability of the hydrogen halides}

The \textbf{thermal stability} 热稳定性 of the hydrogen halides \textbf{decreases} down the group. The H–X bond gets weaker as the halogen atom gets larger, so HI breaks apart on gentle heating while HCl is very stable.

\section*{Reactions of the halide ions}

\subsection*{Halide ions as reducing agents}

A \textbf{halide ion} 卤离子 (such as $\text{Cl}^-$) can give away an electron, acting as a \textbf{reducing agent} 还原剂. This power \textbf{increases} down the group, because a larger ion holds its outer electron less tightly.

\subsection*{Reaction with aqueous silver ions}

Add aqueous silver nitrate, then aqueous \textbf{ammonia} 氨, to identify the halide from the colour of the silver halide \textbf{precipitate} 沉淀:

\begin{center}
\adjustbox{max width=\linewidth}{%
\begin{tabular}{lll}
\toprule
\textbf{Halide} & \textbf{Precipitate with $\text{Ag}^+$} & \textbf{Solubility in ammonia} \\
\midrule
$\text{Cl}^-$ & white & dissolves in dilute ammonia \\
$\text{Br}^-$ & cream & dissolves only in concentrated ammonia \\
$\text{I}^-$ & yellow & insoluble in ammonia \\
\bottomrule
\end{tabular}}
\end{center}

\subsection*{Reaction with concentrated sulfuric acid}

All the halides first give the hydrogen halide. The lower halides are then oxidised by the acid, because they are stronger reducing agents:

\[
\text{NaCl} + \text{H}_2\text{SO}_4 \rightarrow \text{NaHSO}_4 + \text{HCl}
\]

\begin{itemize}
\item chloride gives only $\text{HCl}$ (no redox).

\item bromide also gives some brown $\text{Br}_2$ and $\text{SO}_2$.

\item iodide gives $\text{I}_2$ and the smelly gases $\text{H}_2\text{S}$ and $\text{SO}_2$, because $\text{I}^-$ is the strongest reducing agent.

\end{itemize}

\section*{Reactions of chlorine}

\subsection*{With sodium hydroxide}

With \textbf{cold, dilute} sodium hydroxide, chlorine reacts to form chloride and chlorate(I):

\[
\text{Cl}_2 + 2\text{NaOH} \rightarrow \text{NaCl} + \text{NaClO} + \text{H}_2\text{O}
\]

With \textbf{hot, concentrated} sodium hydroxide, it forms chloride and chlorate(V):

\[
3\text{Cl}_2 + 6\text{NaOH} \rightarrow 5\text{NaCl} + \text{NaClO}_3 + 3\text{H}_2\text{O}
\]

In both, the \textbf{oxidation number} 氧化数 of chlorine goes both up and down (from $0$), so both are \textbf{disproportionation} 歧化 reactions.

\subsection*{Chlorine in water purification}

A little chlorine is added to water for \textbf{water purification} 水净化. It reacts with water:

\[
\text{Cl}_2 + \text{H}_2\text{O} \rightleftharpoons \text{HOCl} + \text{HCl}
\]

The active species $\text{HOCl}$ and $\text{ClO}^-$ kill \textbf{bacteria} 细菌, making the water safe to drink.


\end{document}
